\subsection{Regression checks and the scheduling decision}

All 23 focused tests pass in 3.292 seconds. After the geometry-reuse
refactoring, all 1,210 maintained tests pass in 220.299 seconds. The final
source audit takes 75.389 seconds and the boundary survey takes 57.621
seconds. These are correctness runs, partly overlapping the full test
suite; their elapsed times are not performance comparisons. The final
audit, survey and benchmark each pin 493 source files.

The initial source and survey records, before geometry reuse, are retained.
All surveyed boundary records and all five selected certificates are
identical after refactoring. Verdicts under both source policies are
unchanged, and the off-arm records still exactly reproduce the preceding
annulus audit. The initial timing implementation is preserved as
hash-verified source snapshots, together with the earlier audit driver's
fixture map. This distinguishes the tested revisions from the final
implementation rather than assigning historical measurements its source
hashes.

Both timing runs start from source PDs and include complete recognition,
any exact fallback, independent replay of the returned native certificate,
and certificate hashing. Earlier recognition shortcuts are disabled to
reach the native stage; the seed has no local work cap, and the global
deadline remains twelve seconds. Each case has a warm-up round and five
measured rounds, randomizing off, off A/A, on and on A/A order. Tree trials
are four except for the three new extended positives, which request 24;
face discovery remains disabled. No tests, audits, builds or commits
overlap either timing run.

All 440 measured calls and 88 warm-ups across the initial and revised
implementations complete. Per-implementation verdicts and certificate
hashes agree across both runs. Native work is stable within each revision;
the intentional reductions after geometry reuse are recorded rather than
claimed unchanged. The initial run takes 28.735 seconds and the final run
22.443 seconds. The table reports the final revision only: primary-arm
median times and medians of paired ratios, which need not equal the ratio
of the displayed medians.

\begin{center}
\small\begin{tabular}{lrrrrr}
\toprule
Input & Off (ms) & On (ms) & Off/on & Off A/A & On A/A \\
\midrule
Raw disc control & 61.430 & 64.131 & 0.970 & 1.000 & 0.999 \\
Raw annulus control & 30.756 & 30.544 & 1.007 & 0.989 & 0.999 \\
Earlier seven-boundary proof & 88.652 & 78.822 & 1.121 & 0.995 & 1.006 \\
New four-tree positive & 59.070 & 62.398 & 0.916 & 0.945 & 0.977 \\
Earlier raw planar proof & 88.663 & 89.244 & 0.999 & 1.000 & 0.997 \\
Optimized annulus control & 138.529 & 174.510 & 0.794 & 0.984 & 1.020 \\
New extended positive A & 66.745 & 105.452 & 0.631 & 0.996 & 1.058 \\
New extended positive B & 114.391 & 217.936 & 0.530 & 1.007 & 0.966 \\
New extended positive C & 66.821 & 98.140 & 0.669 & 0.998 & 1.001 \\
Trefoil miss & 41.634 & 45.828 & 0.912 & 0.997 & 0.992 \\
Genus-one miss & 42.727 & 51.690 & 0.830 & 1.002 & 1.003 \\
\bottomrule
\end{tabular}

\end{center}

The earlier seven-boundary witness improves the complete endpoint by a
paired factor of $1.121$: approximately 89 milliseconds becomes 79.
The earlier raw planar witness is near parity at 89 milliseconds.
The new four-tree positive is modestly slower, with a 0.916 off/on ratio
and noticeable off-arm A/A variation. The three new extended positives
are materially slower: their off/on ratios are 0.631, 0.530 and 0.669.
The existing optimized-annulus positive also slows because earlier
unsuccessful candidates now receive boundary queries; its ratio is 0.794.
The trefoil and genus-one continuing misses have ratios 0.912 and 0.831.
Raw disc and raw annulus controls bypass boundary discovery and show no
algorithmic speedup.

These data support an optional stronger certificate search, not a new
universal default. Geometry reuse removes much of the initial overhead,
but the extra queries and larger certificates can still cost more than
the preceding exact fallback. The maintained default therefore remains
\code{planar=False}. The whole native stage remains optional in the
ordinary portfolio, whose nineteen measured inputs never reach it.

Reproduction uses \path{../fast/normal_orbit_research/planar.py}.
Its modes are \code{audit}, \code{survey} and
\code{recognize --rounds 5}. Regina and expanded boundary enumeration
appear only in the audit controls. The incoming research inventory and
checked \code{openai/math} head \code{fd4aeeb2e} are unchanged. The new
capping proof uses elementary surface and torus topology and does not
require a zero-surgery theorem. No complete genus-zero search, explicit
normal capped-disc extraction, hierarchy cutting, or general
subexponential or quasipolynomial recognition bound is established here.
